% !PRINTED 483
\weberpage{483}
\begin{center}
\S.\,151.

\gesperrt{Primitive and imprimitive groups.}
\end{center}

Now let $f(x) = 0$ be an irreducible equation of degree $n$, so that its
\gesperrt{Galois} group $P$ is transitive. Let the roots of $f(x)$ be
\weq{1}{\alpha, \alpha_1, \alpha_2 \ldots \alpha_{n-1}.}

If the field $\Omega(\alpha)$ is imprimitive (\S.\,144) and
$\Theta = \chi(\alpha)$ is an imprimitive element whose conjugate values fall
into $s$ systems of $r$ mutually equal values each, so that
\[
n = rs
\]
with $r > 1$, $s < n$, then the values (1) can be decomposed into $s$
sequences of $r$ elements each, which we shall denote by
\weq{2}{
\begin{array}{l}
A = \alpha, \alpha_1 \ldots \alpha_{r-1} \\
B = \beta, \beta_1 \ldots \beta_{r-1} \\
\cdots\cdots\cdots\cdots \\
S = \sigma, \sigma_1 \ldots \sigma_{r-1}
\end{array}
}
so that
\weq{3}{
\begin{array}{l}
\Theta = \chi(\alpha) = \chi(\alpha_1) \ldots = \chi(\alpha_{r-1}) \\
\Theta_1 = \chi(\beta) = \chi(\beta_1) \ldots = \chi(\beta_{r-1}) \\
\cdots\cdots\cdots\cdots\cdots\cdots\cdots\cdots \\
\Theta_{s-1} = \chi(\sigma) = \chi(\sigma_1) \ldots = \chi(\sigma_{r-1})
\end{array}
}
are the conjugate values of $\Theta$. By \S.\,144,
\weq{4}{(t - \Theta)(t - \Theta_1) \ldots (t - \Theta_{s-1}) = \varphi(t)}
is then an irreducible function in $\Omega$ of degree $s$ with respect to $t$,
whose roots are the values (3).

It now follows from (3) that the group $P$ must be so constituted that
all its permutations interchange the elements of the individual sequences
$A, B \ldots S$ only among themselves, and moreover permute the whole
sequences $A, B \ldots S$ with one another, so that two elements of one
sequence are never replaced by two elements of different sequences. For if,
say, a permutation $\pi$ of $P$ carried $\alpha$ and $\alpha_1$ into $\beta$
and $\sigma$, it would follow, since one may perform the permutation $\pi$
[by a), \S.\,149] in the equation $\chi(\alpha) = \chi(\alpha_1)$, that
$\chi(\beta) = \chi(\sigma)$ would also have to hold, which contra-
%CONT>>
